%\!TEX program = xelatex
\documentclass[12pt,a4paper]{report}
\usepackage[utf-8]{inputenc}
\usepackage{geometry}
\geometry{margin=1in}
\usepackage{graphicx}
\usepackage{hyperref}
\usepackage{amsmath}
\usepackage{algorithm}
\usepackage{algpseudocode}
\usepackage{cite}
\usepackage{fancyhdr}
\usepackage{setspace}
\onehalfspacing

\title{Hybrid Explainable AI Framework for Deep Learning Models}
\author{Jankit Shah \\ University of East London}
\date{\today}

\begin{document}

\maketitle

\begin{abstract}
This dissertation presents a comprehensive hybrid framework that combines three state-of-the-art explainability methods—LIME, SHAP, and Grad-CAM—to provide interpretable explanations for deep learning models. The framework is designed to be model-agnostic and domain-flexible, supporting multiple deep learning architectures (CNNs, RNNs, Transformers) across different application domains (Computer Vision, NLP, Time Series). We evaluate the framework's effectiveness through comparative analysis, demonstrating that combining multiple explanation methods provides more robust and trustworthy interpretability. Our findings suggest that a hybrid approach overcomes limitations of individual methods and provides practitioners with a powerful tool for understanding and debugging deep learning models.
\end{abstract}

\tableofcontents

\chapter{Introduction}

\section{Background}
Deep learning models have achieved remarkable success across various domains, but their complex nature has made them difficult to interpret. This ``black box'' nature poses significant challenges in high-stakes applications such as healthcare, finance, and autonomous systems, where understanding model decisions is crucial.

\section{Motivation}
Explainability is not just an academic concern—it has become a regulatory and ethical imperative. GDPR's "right to explanation" and similar regulations require organizations to explain algorithmic decisions. Furthermore, practitioners need tools to debug models, identify biases, and build trust in AI systems.

\section{Objectives}
This dissertation aims to:
\begin{enumerate}
    \item Design and implement a hybrid explainability framework
    \item Integrate LIME, SHAP, and Grad-CAM in a unified interface
    \item Support multiple model architectures and domains
    \item Evaluate and compare explanation methods
    \item Provide practical tools for practitioners
\end{enumerate}

\chapter{Literature Review}

\section{Explainable AI}
Explainable AI (XAI) is a research field focused on creating machine learning models whose actions can be understood by domain experts \cite{samek2017explainable}.

\section{LIME}
LIME (Local Interpretable Model-agnostic Explanations) is a technique to explain predictions of any classifier by approximating it locally with an interpretable model.

\section{SHAP}
SHAP (SHapley Additive exPlanations) values are a unified measure of feature importance based on cooperative game theory.

\section{Grad-CAM}
Gradient-weighted Class Activation Mapping (Grad-CAM) uses the gradients flowing into the final convolutional layer to produce a localization map of important regions.

\chapter{Framework Design}

\section{Architecture}
\subsection{Core Components}
\begin{itemize}
    \item Model Wrapper: Abstracts different model types
    \item Explainer Interface: Unified API for all methods
    \item Visualization Engine: Comparative and individual visualizations
    \item Evaluation Module: Metrics for explanation quality
\end{itemize}

\section{Implementation Details}
The framework is implemented in Python using PyTorch for deep learning, with LIME, SHAP, and Captum libraries for the explanation methods.

\chapter{Experimental Evaluation}

\section{Experimental Setup}
\subsection{Datasets}
\begin{itemize}
    \item Computer Vision: CIFAR-10, ImageNet subset
    \item NLP: Movie Review Dataset
    \item Time Series: Stock price prediction
\end{itemize}

\section{Results}
\subsection{Comparative Analysis}
Our results demonstrate that different explanation methods provide complementary insights into model behavior.

\chapter{Conclusions and Future Work}

\section{Conclusions}
The hybrid explainability framework successfully combines multiple explanation methods, providing practitioners with a comprehensive toolkit for understanding deep learning models.

\section{Future Work}
Potential directions for future work include:
\begin{itemize}
    \item Support for additional explanation methods
    \item Integration with model debugging tools
    \item Interactive visualization improvements
    \item Automated explanation generation
\end{itemize}

\bibliographystyle{plain}
\begin{thebibliography}{99}

\bibitem{samek2017explainable} Samek, W., Montavon, G., & Kawanabe, M. (2017). Explainable artificial intelligence. In Explainable AI (pp. 1-3).

\end{thebibliography}

\end{document}
